
\documentclass[9pt]{elife}
\usepackage{amsthm}
\usepackage{amsmath,amssymb,mathtools}
\usepackage{siunitx}
\usepackage{graphicx}
\usepackage{booktabs}
\usepackage{hyperref}
\usepackage{tcolorbox}
\usepackage{bm}

\title{MacroIR: Analytically exact interval likelihoods for time-averaged biological signals}

\author[1*]{Luciano Moffatt}
\affil[1]{Instituto de Qu\'imica F\'isica de los Materiales, Medio Ambiente y Energ\'ia (INQUIMAE), CONICET; Facultad de Ciencias Exactas y Naturales, Universidad de Buenos Aires, Ciudad de Buenos Aires, C1428EHA, Argentina}
\corr{lmoffatt@qi.fcen.uba.ar}

\begin{document}
\maketitle

\begin{abstract}
Time-averaged measurements dominate experimental biology, yet most inference methods assume instantaneous observations, yielding biased estimates when intrinsic dynamics outpace acquisition. MacroIR (Macroscopic Interval Recursive) provides an analytically exact formulation of the likelihood for interval-averaged signals generated by continuous-time Markov processes. By conditioning on both the start and end states of each sampling window, MacroIR computes the interval mean and variance in closed form and performs recursive, Kalman-like updates in the original state space, avoiding the exponential enumeration of trajectories and reducing the quadratic blow-up of explicit start–end  combination states. Within a linear-Gaussian observation model, this yields unbiased dynamic inference and an internal diagnostic—the score-based self-consistency test—that confirms informational consistency between data sampling and likelihood evaluation. Applications to simulated and experimental data demonstrate that MacroIR corrects systematic bias in time-averaged biological recordings and enables reliable estimation of underlying dynamics across molecular, cellular, and behavioral timescales. 
\end{abstract}

% Impact Statement (for cover letter / submission form):
% MacroIR provides exact, tractable likelihoods for interval-averaged biological data, eliminating long-standing bias and enabling reliable kinetic inference across electrophysiology, single-molecule, and population studies.

\section*{Introduction}

Time-averaged measurements are ubiquitous in experimental biology—electrophysiological amplifiers integrate currents over microseconds, fluorescence detectors average photon flux over milliseconds, and population assays aggregate signals over hours \citep{Hille2001, Denk1990, Elowitz2002}. 
Yet most inference methods treat such data as instantaneous samples. When intrinsic dynamics evolve on comparable timescales, this mismatch between physical averaging and model assumptions distorts likelihoods and biases inferred parameters \citep{ColquhounHawkes1982, Qin2000, Munsky2009, Moffatt}.

Continuous-time Markov models (CTMCs) remain the standard framework for describing stochastic molecular and cellular dynamics. Over recent decades, inference methods have advanced from rate-equation and hidden-Markov formulations to recursive Bayesian and Kalman-filter approaches \citep{Csanady2006, Siekmann2011, Bronson2013}. 
Münch \emph{et al.} introduced a generalized Bayesian Kalman filter for ion-channel kinetics that unifies current and fluorescence data while incorporating state-dependent open-channel noise and finite temporal resolution \citep{Munch2022}. 
Their formulation integrates mean dynamics exactly over each sampling window through the propagator $T = e^{K\Delta t}$ and propagates the process covariance $Q(T,\bar n)$ by integrating infinitesimal fluctuations under a Gaussian occupancy approximation. 
As a result, the posterior ensemble state effectively represents the \emph{average occupancy during the interval} in the linear–Gaussian regime, providing accurate inference when the underlying dynamics vary slowly relative to the sampling step.

However, in this implicit treatment of interval averaging, when acquisition windows approach the system’s dwell times, the Gaussian approximation breaks down and structural bias arises in rate estimation and model evidence.

Here we introduce \textbf{MacroIR} (Macroscopic Interval Recursive), a framework that makes interval averaging explicit within CTMC-based inference. 
MacroIR derives closed-form expressions for the mean and variance of \emph{interval-averaged observables} conditioned on both the start and end states of each measurement window and performs recursive, Kalman-like updates directly in the original state space. 
Like the approach of Münch \emph{et al.}, it avoids exponential trajectory enumeration and quadratic scaling with the number of states. 
Within a linear–Gaussian observation model, MacroIR yields unbiased inference of dynamic parameters.

We also introduce a diagnostic—the \emph{score-based self-consistency test}—that quantifies informational agreement between sampling and likelihood. 
We demonstrate its application to macroscopic current recordings and discuss its generalization across molecular, cellular, and population-level systems.


\section*{Theory}

\subsection*{Interval-averaged observation model}

Let $x_t$ denote a continuous-time Markov chain (CTMC) with generator $Q$ and propagator $P(t)=e^{Qt}$.  
For a linear observable $\bm{\gamma}\in\mathbb{R}^k$, the measurement obtained over an acquisition window of duration $t$ is the interval average
\[
\bar y_{0\to t} = \frac{1}{t}\int_0^t \gamma(x_s)\,ds.
\]
Conditioned on the initial and final states $(i,j)$, the expected interval average is
\begin{equation}
	\gamma_{i\to j}(t)
	= \frac{1}{P_{i\to j}(t)}
	\sum_{n_1,n_2}
	V_{i n_1}\,(V^{-1}\Gamma V)_{n_1 n_2}\,V^{-1}_{n_2 j}\,
	E_2(\lambda_{n_1} t, \lambda_{n_2} t),
\end{equation}
where $Q=V\Lambda V^{-1}$, $\Gamma=\mathrm{diag}(\bm{\gamma})$, and
\[
E_2(x,y) = \frac{e^{x}-e^{y}}{x-y}
\]
is the two-exponential kernel that arises from integrating the propagator over the interval.

\begin{tcolorbox}[colback=gray!5,colframe=gray!60,title=Boundary-conditioned two-point states (“meta-states”)]
	Define the ordered pair $(i,j)$ representing start and end states over duration $t$.  
	The $k^2$ possible pairs form a space whose generator is the Kronecker sum
	$Q\oplus Q = Q\otimes I + I\otimes Q$.  
	Because the propagator factorizes as $e^{(Q\oplus Q)t}=e^{Qt}\otimes e^{Qt}$, all required expectations reduce to bilinear forms in the original $k$-dimensional state space—no explicit $k^2$ object is needed.
\end{tcolorbox}

\subsection*{Recursive update and gain}

Given a prior state distribution $(\bm{\mu}^{\mathrm{prior}}_0,\Sigma^{\mathrm{prior}}_0)$ and an observed interval-averaged measurement $y^{\mathrm{obs}}_{0\to t}$, the predicted mean and variance of the observable are
\begin{align}
	y^{\mathrm{pred}}_{0\to t} &= N_{\mathrm{ch}}\,(\bm{\mu}^{\mathrm{prior}}_0\!\cdot\!\bm{\gamma}_0),\\[3pt]
	\sigma^2_{y^{\mathrm{pred}}_{0\to t}} &= 
	\frac{\epsilon^2}{t} + \nu^2 
	+ N_{\mathrm{ch}}\!\left[\bm{\gamma}_0^\top(\Sigma^{\mathrm{prior}}_0-\mathrm{diag}\,\bm{\mu}^{\mathrm{prior}}_0)\bm{\gamma}_0 
	+ \bm{\mu}^{\mathrm{prior}}_0\!\cdot\!\mathbf{E}\right],
\end{align}
where $\bm{\gamma}_0$ is the vector of \emph{interval-averaged conductances} defined in Eq.~(1), i.e. $\bm{\gamma}_0=\{\gamma_{i\to j}(t)\}$, which depends on the sampling duration $t$ and is conditioned on the start state.  
The total observation noise consists of three additive components:
a photon (shot) noise term $\epsilon^2/t$,
a white instrumental noise term $\nu^2_{\mathrm{white}}$, 
and a pink (1/f) noise term $\nu^2_{\mathrm{pink}}$ that remains constant with $t$, representing low-frequency baseline drift not captured by the Markov dynamics. 
Hence, the effective observation variance becomes
\[
\sigma^2_{y^{\mathrm{pred}}_{0\to t}} = 
\frac{\epsilon^2}{t} + \nu^2_{\mathrm{white}} + \nu^2_{\mathrm{pink}}
+ N_{\mathrm{ch}}\!\left[\bm{\gamma}_0^\top(\Sigma^{\mathrm{prior}}_0-\mathrm{diag}\,\bm{\mu}^{\mathrm{prior}}_0)\bm{\gamma}_0 
+ \bm{\mu}^{\mathrm{prior}}_0\!\cdot\!\mathbf{E}\right].
\]


The posterior moments are updated in a Kalman-like form:
\begin{align}
	\bm{\mu}^{\mathrm{post}}_{t} &=
	\bm{\mu}^{\mathrm{prior}}_0 P(t)
	+ \frac{y^{\mathrm{obs}}_{0\to t}-y^{\mathrm{pred}}_{0\to t}}{\sigma^2_{y^{\mathrm{pred}}_{0\to t}}}\,K^\mu_{0\to t},\\[3pt]
	\Sigma^{\mathrm{post}}_{t} &=
	\Sigma^{\mathrm{prior}}_{t}
	- \frac{1}{\sigma^2_{y^{\mathrm{pred}}_{0\to t}}}
	(K^\mu_{0\to t})^\top K^\mu_{0\to t},
\end{align}
with Kalman gain
\[
K^\mu_{0\to t} = \frac{\mathrm{Cov}(x_t,\bar y_{0\to t})}{\mathrm{Var}(\bar y_{0\to t})},
\]
which admits a closed-form expression from $P(t)$, $\gamma_{i\to j}(t)$, and prior moments.

% \begin{figure}[t]
	% \centering
	% \includegraphics[width=0.85\linewidth]{fig2_theory.pdf}
	% \caption{\textbf{Theory schematic.} (A) Boundary conditioning. (B) Bilinear factorization. (C) Kalman-like gain mapping interval residuals to state updates.}
	% \label{fig:theory}
	% \end{figure}
\section*{Methods}
\subsection*{Synthetic models and data}
We simulate $k$-state CTMCs with known $Q$ and linear observables $\bm{\gamma}$; traces are averaged over bins $\Delta$ spanning sub- to supra-kinetic regimes. White/pink noise $(\epsilon^2,\nu^2)$ follow baseline estimates. Priors on rates are log-uniform; noise parameters use weakly-informative priors.

\subsection*{Comparative benchmarks}
We implement three alternatives: (i) midpoint likelihood (single-point per bin), (ii) non-recursive interval likelihood (NR) without boundary conditioning, and (iii) instantaneous-sample approximation. For each method we report bias, root-mean-square error (RMSE), coverage of 95\% intervals, and runtime. MacroIR uses identical priors and samplers. Evidence via thermodynamic integration is used when model ranking is evaluated.

\subsection*{Score test}
Let $S(\theta)=\nabla_\theta \log \mathcal{L}(\theta)$ be the per-interval score. Under correct likelihood, $\mathbb{E}[S]=0$ and $\mathrm{Var}(S)=\mathcal{I}(\theta)$ (Fisher). We estimate mean/variance of $S$ and compare to a plug-in $\widehat{\mathcal{I}}$ from MacroIR. Diagnostics are shown across $\Delta/\tau_{\min}$ and SNR.

% \begin{table}[t]
% \centering
% \caption{\textbf{Benchmark summary.} Mean bias, RMSE, coverage, and runtime across methods.}
% \begin{tabular}{lcccc}
% \toprule
% Method & Bias & RMSE & 95\% coverage & Runtime \\
% \midrule
% MacroIR & & & & \\
% Midpoint & & & & \\
% NR & & & & \\
% Instantaneous & & & & \\
% \bottomrule
% \end{tabular}
% \label{tab:bench}
% \end{table}

\section*{Results}
\subsection*{Synthetic recovery and benchmark comparisons}
MacroIR recovers unbiased parameters across $\Delta/\tau_{\min}$, with calibrated uncertainty and higher evidence for true models. Midpoint and NR display increasing bias and under-coverage as $\Delta$ grows, and instantaneous approximations fail when $\Delta\gtrsim \tau_{\min}$.

% \begin{figure}[t]
% \centering
% \includegraphics[width=0.85\linewidth]{fig3_benchmark.pdf}
% \caption{\textbf{Synthetic benchmarks.} (A) Bias vs $\Delta/\tau_{\min}$. (B) Coverage vs $\Delta/\tau_{\min}$. (C) Log-evidence difference vs ground truth.}
% \label{fig:bench}
% \end{figure}

\subsection*{Score-based self-consistency}
Only MacroIR yields mean score $\approx 0$ with variance matching Fisher across regimes; alternatives show nonzero means and underestimation of information, especially at large $\Delta$.

% \begin{figure}[t]
% \centering
% \includegraphics[width=0.8\linewidth]{fig4_score.pdf}
% \caption{\textbf{Score diagnostics.} (A) Mean score vs $\Delta/\tau_{\min}$. (B) $\mathrm{Var}(S)/\mathcal{I}$ ratio across methods.}
% \label{fig:score}
% \end{figure}

\subsection*{Biological case study: P2X receptor kinetics}
Reanalysis of ultrafast ATP-pulse macroscopic currents from P2X$_2$ receptors shows MacroIR reproduces responses and correctly ranks candidate schemes, revealing asymmetric coupling between ligand-binding and subunit rotation. This mechanistic asymmetry was not supported by competing approximations, aligning with the biological interpretation detailed in Moffatt (2025, \emph{Communications Biology}, in press).

% \begin{figure}[t]
% \centering
% \includegraphics[width=0.85\linewidth]{fig5_p2x.pdf}
% \caption{\textbf{P2X case study.} (A) Data and MacroIR fit. (B) Scheme ranking (evidence). (C) Posterior on asymmetry parameters.}
% \label{fig:p2x}
% \end{figure}

\subsection*{Computational performance}
Cost scales linearly with the number of intervals and as $O(k^3)$ with state dimension, dominated by matrix exponentials and linear solves. Including an explicit analog filter via state-space augmentation increases runtime $\sim$10$\times$ but remains tractable for moderate $k$.

% \begin{figure}[t]
% \centering
% \includegraphics[width=0.78\linewidth]{fig6_scaling.pdf}
% \caption{\textbf{Scaling.} Wall-clock time vs $k$ and intervals; dashed line shows $O(k^3)$.}
% \label{fig:scaling}
% \end{figure}

\section*{Discussion}
MacroIR restores informational consistency between sampling and inference for time-averaged observations, eliminating bias endemic to instantaneous-likelihood assumptions. The framework is exact for true integrals and offers a built-in diagnostic via the score test. It applies broadly beyond ion channels to fluorescence trajectories, gene expression bursts, and population dynamics.

\subsection*{Limitations and outlook}
(1) Scaling may become demanding for very large $k$; sparsity/exponential Krylov methods could help. (2) Explicit analog filtering requires state-space augmentation or kernel convolution and increases runtime. (3) Robustness to model misspecification deserves study; extensions to mixture or hierarchical CTMCs are natural. (4) Nonlinear observables would require linearization or alternative moment closures.

\section*{Data and Code Availability}
Code and scripts to reproduce all results will be released publicly (GitHub/Zenodo) at submission, including minimal working examples and figure-generation notebooks.

\section*{Acknowledgements}
I thank colleagues for discussions on boundary conditioning, and local HPC facilities for computational support.

\section*{Author Contributions}
LM conceived the study, developed the theory, implemented algorithms, analyzed data, and wrote the manuscript.

%\nocite{*} % This command displays all refs in the bib file. PLEASE DELETE IT BEFORE YOU SUBMIT YOUR MANUSCRIPT!

\bibliography{biblio}

\end{document}
